\section{System Design}

OriginBlame stores all provenance metadata in a \texttt{.ob/} directory colocated with project data as plain JSONL---no database, no server, no configuration files.

\subsection{Architecture}

OriginBlame employs a three-tier content-addressable model (Figure~\ref{fig:three-layer})---each entity is identified by the SHA-256 hash~\cite{nist2015fips1804} of its contents, not an assigned ID. A key distinction separates \emph{authors} (editors whose content is currently visible) from \emph{contributors} (prior editors whose changes were overwritten). Authors hold revocation rights; contributors are tracked for audit but have no revocation authority. (ob stores contributor ids but does not expose \texttt{revoke --contributor}, since overwritten edits have no content to withdraw.)

The \textbf{authors layer} stores identities (\texttt{id}, \texttt{name}, \texttt{email}, \texttt{revoked}). The \textbf{sections layer} stores copyright registration entries: each section has a \texttt{section\_hash} (SHA-256 of \texttt{\{path, authors, contributors, license, year\}}), \texttt{path}, \texttt{license}, \texttt{year}, and \texttt{revoked}. The \textbf{document-index layer} links each output record to its sources via \texttt{line\_hash}, \texttt{file}, and \texttt{sources} (section hashes). All files are sharded into 256 buckets by hash prefix for expected O(1) lookup; multi-process safety uses process-ID-isolated writes merged by \texttt{ob clean}.

\subsection{Three-Level Revocation}

Revocation operates at three granularities, all via a boolean tag toggle in O(1) time with lazy cascade at query time (tags propagate on read rather than being eagerly propagated on write, so revocation itself is instantaneous regardless of dataset size):

\noindent\textbf{Author-level} (\texttt{ob revoke -{}-author}): sets \texttt{revoked=true} on the author record. Queries lazily cascade: authors $\rightarrow$ sections referencing that author $\rightarrow$ document-index entries referencing those sections. The result is a precise \texttt{(file, line\_hash)} list of every data line traceable to the revoked author.

\noindent\textbf{Section-level} (\texttt{ob revoke -{}-section HASH}): toggles a single section's tag, withdrawing one file or data source without affecting other content from the same author.

\noindent\textbf{Record-level} (\texttt{ob revoke -{}-line-hash HASH -{}-file FILE}): toggles a single document-index entry, enabling finest-grain withdrawal of individual records.

All revocations are reversible via \texttt{-{}-reverse}. \texttt{ob purge} physically deletes revoked lines from data files; a safety constraint enforces no purge without prior revoke.

\subsection{Token-Level Provenance}

Subword tokenizers split records at arbitrary boundaries, and packing concatenates multiple records into fixed-length training sequences---together destroying per-record identity. ob addresses this with an independent \textbf{token-index} layer that links each document's token range in the packed binary to its source chain. Each entry stores \texttt{token\_count}, \texttt{sources} (section hashes), \texttt{tokenizer} (e.g., ``gpt2''~\cite{radford2019gpt2}), and \texttt{revoked}. The token-index operates in two modes: \emph{file mode} (coexisting with document-index) and \emph{streaming mode} (tokenizing directly without intermediate JSONL). \texttt{ob generate-set} produces a binary bitmask directly usable by unlearning algorithms.

\subsection{Pipeline Integration}

Integration follows a four-step one-time setup: \texttt{ob init}, \texttt{ob author.add}, \texttt{ob register.add}, and \texttt{source.append(path)}. Once configured, the pipeline inserts a single \texttt{track(record, file)} call where each training record is written. Internally, \texttt{track} computes the SHA-256 hash, resolves active sources, checks for duplicates, and writes to process-ID-isolated files under write-ahead logging (WAL) for crash recovery. ob is a non-invasive observer---no pipeline logic changes required.

We validate framework-agnostic integration with two pipelines. For Datatrove~\cite{penedo2024datatrove}, an \texttt{OBProvenanceBlock} PipelineStep reads contributor metadata from \texttt{Document.metadata} and writes token-index entries. For HuggingFace Datasets~\cite{lhoest2021datasets}, a \texttt{.map()} callback wraps \texttt{track()} around each record. Both adapters require $\sim$5 lines and produce identical \texttt{.ob/} directories; experimental overhead is reported in Section~\ref{sec:evaluation}.
